\documentclass[french]{article}
\usepackage[utf8]{inputenc}
\usepackage[T1]{fontenc}
\usepackage{lmodern}
\usepackage{geometry}
\usepackage{graphicx}
\usepackage{babel}
\usepackage{amsmath}
\usepackage{amsthm}
\usepackage{amssymb}
\usepackage{hyperref}
\usepackage{stmaryrd}
\usepackage{enumitem}
\usepackage{tcolorbox}
\usepackage{dsfont}
\usepackage{multirow}
\usepackage{etoc}
\usepackage{titlesec}
\usepackage{esint}
\usepackage{cancel}
\usepackage{mathdots}

\setlist[itemize]{label=$\bullet$}


\renewcommand{\P}{\mathbb{P}}
\newcommand{\N}{\mathbb{N}}
\newcommand{\Z}{\mathbb{Z}}
\newcommand{\Q}{\mathbb{Q}}
\newcommand{\R}{\mathbb{R}}
\newcommand{\C}{\mathbb{C}}
\newcommand{\K}{\mathbb{K}}
\renewcommand{\L}{\mathbb{L}}
\newcommand{\U}{\mathbb{U}}
\newcommand{\st}{^*}
\newcommand{\tq}{\middle|}
\newcommand{\inv}{^{-1}}
\newcommand{\E}{\mathbb{E}}
\newcommand{\F}{\mathbb{F}}
\newcommand{\V}{\mathbb{V}}
\renewcommand{\ker}{\text{Ker}}
\newcommand{\im}{\text{Im}}
\newcommand{\card}{\text{Card}}
\newcommand{\Tr}{\text{Tr}}

\title{TD Euclidiens\\ Exercice 34}
\date{}
\begin{document}
\maketitle

\section{Énoncé}
Soit $u$ un endomorphisme d'un espace euclidien $E$. Montrer l'équivalence des propriétés suivantes : 
\begin{enumerate}
	\item $u$ commute avec $u\st$ ;
	\item $\forall x\in E,\ \lVert u(x)\rVert=\lVert u\st(x)\rVert$ ;
	\item $\forall x\in E,\ \lVert u(x)\rVert\leqslant \lVert u\st(x)\rVert$.
\end{enumerate}
On appelle endomorphisme normal tout endomorphisme vérifiant ces trois conditions équivalentes.
\section{Solution}
On procède à des implications circulaires.
\begin{itemize}
	\item $1\implies 2$ : soit $x\in E$. On a : $$\lVert u(x)\rVert2=(u(x)\mid u(x))=(x\mid u\st\circ u(x))=(x\mid u\circ u\st(x))=(u\st(x)\mid u\st(x))=\lVert u\st(x)\rVert^2$$ donc par positivité, $\lVert u(x)\rVert=\lVert u\st(x)\rVert$
	\item $2\implies 3$ : immédiat
	\item $3\implies 1$ : alors pour tout $x\in E$, par bilinéarité et symétrie : $$\left(\underbrace{u\st \circ u-u\circ u\st}_{=v}(x)\mid x\right)\leqslant 0$$ Or, il est clair que $v\st=v$ donc on en déduit que $(v(x)\mid x)\geqslant0$ et donc $(v(x)\mid x)=0$. Puis comme $v$ est diagonalisable, toutes ses valeurs propres sont nulles et $v=0$ (se référer à un exercice de la partie précédente pour les détails). Donc $u$ et $u\st$ commutent.
\end{itemize}
\end{document}
